% mage-auto-heading: component-and-zone
% mage-auto-heading: service-flow
% mage-auto-heading: two-simpler-structural-forms
% Authoritative LaTeX manuscript; Markdown is a generated compatibility projection.

Structural models represent what exists and how elements relate. In review, they expose which architectural boundaries a change affects. Ownership and boundary policy often share the same representation but answer distinct questions: who owns a surface, and which relationships are permitted.

\section{Component and Zone}\label{heading:component-and-zone}

\textbf{Engineering question.} Where am I, who owns this surface, and what may cross its boundary?

The model maps implementation surfaces to components and declares the seams through which those components may interact.

\textbf{Representation.} A minimal form:

% mage-code-language: text
\begin{verbatim}
Component
    id
    name
    focus_directories[]
    tags[]
    boundary_kind
Zone
    component_id
    directory
Seam
    source_component
    target_surface
    relation_kind
    permitted
\end{verbatim}

\autoref{fig-g2-structure-boundaries} joins the authored catalog to the observed tree.

\begin{figure}[htbp]
\centering
\includegraphics[width=\linewidth]{book/latex/assets/appendix-g-2-structure-boundaries-5eea5569e4.pdf}
\caption{\emph{Structure and boundaries.} Authored ownership and permitted seams are reconciled with the observed repository tree. Sanctioned crossings pass through declared doors; forbidden direct access raises a finding.}
\label{fig-g2-structure-boundaries}
\end{figure}

\textbf{Property.} Representative properties, each with the source that decides it:

{\footnotesize
\begin{longtable}[]{@{}
  >{\raggedright\arraybackslash}p{(\linewidth - 2\tabcolsep) * \real{0.5000}}
  >{\raggedright\arraybackslash}p{(\linewidth - 2\tabcolsep) * \real{0.5000}}@{}}
\caption[Component-zone properties]{Representative properties of the component-and-zone model, and the fact source each draws on.}\tabularnewline
\toprule\noalign{}
\begin{minipage}[b]{\linewidth}\raggedright
Property
\end{minipage} & \begin{minipage}[b]{\linewidth}\raggedright
Fact source
\end{minipage} \\
\midrule\noalign{}
\endfirsthead
\toprule\noalign{}
\begin{minipage}[b]{\linewidth}\raggedright
Property
\end{minipage} & \begin{minipage}[b]{\linewidth}\raggedright
Fact source
\end{minipage} \\
\midrule\noalign{}
\endhead
\bottomrule\noalign{}
\endlastfoot
Every in-scope source directory belongs to exactly one component & Authored ownership joined to the observed tree \\
Every component\textquotesingle s declared directory exists & Repository \\
Boundary crossings use sanctioned seams & Authored model \\
Powerful mutation surfaces are reachable only through declared doors & Authored model plus observed dependency \\
\end{longtable}
}

Directory existence can be derived from the repository; ownership and permitted boundaries are authored architectural decisions.

\textbf{Enforcement and correspondence.} The repository tree is mapped back to the component catalog. A directory with no owner, multiple owners, or an observed dependency crossing an undeclared boundary becomes a correspondence finding. Where reconciliation depends on source locations, anchor references to resolvable symbols rather than line numbers so the correspondence can be re-derived after edits. The finding may remain advisory or feed a lint, build failure, or admission gate.

\section{Service Flow}\label{heading:service-flow}

\textbf{Engineering question.} Which service may call or reach which service or resource, and under what policy?

\textbf{Representation.}

% mage-code-language: text
\begin{verbatim}
Service
    id
    role
Resource
    id
    kind
Flow
    source
    target
    relation
    auth_posture
\end{verbatim}

The central object is the declared edge. Each permitted service-to-service or service-to-resource relation records both reachability and its authentication posture; an observed edge with no declaration is a mismatch.

\textbf{Property.} Only declared service-to-service and service-to-resource relationships are permitted, each with its declared authentication posture. The same graph carries both connectivity and access policy.

\textbf{Enforcement and correspondence.} The declared graph supplies the \emph{ought}. Static call sites, generated configuration, deployment wiring, and runtime observation each supply part of the \emph{is}. A correspondence check asks whether each observed edge is declared and, where it matters, whether each declared edge still exists.

The same representation pattern applies to firewall policy, service-mesh authorization, network policy, and cloud access control: the representation declares permitted relations while the enforcement mechanism varies by substrate.

\section{Two Simpler Structural Forms}\label{heading:two-simpler-structural-forms}

Two simpler forms use the same pattern: reconcile observed facts or edges against a declared set.

\begin{itemize}
\tightlist
\item
  \textbf{Domain registry.} A registry stores a slowly changing fact under one stable key. Consumers query or join through that key rather than maintaining independent copies.
\item
  \textbf{Bill of materials.} A bill of materials records the third-party dependencies present in a build. Where an authored dependency set exists, completeness can be checked by reconciling the observed dependency edges against it.
\end{itemize}
